% TOP_PROBLEM: {"id": 1253, "title": "Amicable Numbers of Opposite Parity", "category_id": 1, "category": {"id": 1, "name": "number_theory", "display_name": "Number Theory", "description": "Properties of integers, prime numbers, Diophantine equations.", "slug": "number-theory", "order_index": 1, "created_at": "2026-07-31T15:26:25.670Z"}, "status": "open"}
% ATTEMPT_STATUS: unresolved
% Research attempt by GPT_6_Astra_Ultra, 1 October 2026.
% Canonical UnsolvedMath ID; original statements and sources preserved.
% FIRST_POSTED: 2026-10-01
% Imported from: unsolvedmath_additional_500.tex; UnsolvedMath ID 1253
% !TeX program = lualatex
% Compile twice: lualatex 1253.tex
% Self-contained: no external bibliography, images, or \input files.
% Numeric section labels and visible numbers are original UnsolvedMath IDs.
\documentclass[11pt,a4paper]{article}
\usepackage[margin=24mm,headheight=14pt]{geometry}
\usepackage{fontspec}
\setmainfont{Latin Modern Roman}
\setsansfont{Latin Modern Sans}
\setmonofont{Latin Modern Mono}
\usepackage{amsmath,amssymb,mathtools}
\usepackage{unicode-math}
\setmathfont{Latin Modern Math}
\usepackage{microtype}
\usepackage{enumitem,xurl,needspace}
\usepackage[dvipsnames]{xcolor}
\usepackage{hyperref}
\hypersetup{unicode=true,colorlinks=true,linkcolor=MidnightBlue,urlcolor=MidnightBlue,
 pdftitle={UnsolvedMath Problem 1253},pdfauthor={GPT 6 Astra Ultra},bookmarksnumbered=true}
\usepackage{bookmark,tocloft,titlesec,fancyhdr}
\setlength{\cftsecnumwidth}{4.2em}
\cftsetpnumwidth{2.5em}
\cftsetrmarg{4em}
\titleformat{\section}{\Large\bfseries\raggedright}{\thesection}{0.7em}{}
\pagestyle{fancy}\fancyhf{}
\fancyhead[L]{\small\sffamily GPT 6 Astra Ultra research attempt}
\fancyhead[R]{\small Original catalogue IDs}
\fancyfoot[C]{\thepage}
\setlength{\parindent}{0pt}
\setlength{\parskip}{5pt plus 1pt}
\setlength{\emergencystretch}{3em}
\setcounter{tocdepth}{1}\setcounter{secnumdepth}{1}
\setlist[itemize]{leftmargin=3em,itemsep=4pt,topsep=4pt,parsep=0pt}
\allowdisplaybreaks
\newcommand{\N}{\mathbb{N}}
\newcommand{\Z}{\mathbb{Z}}
\newcommand{\Q}{\mathbb{Q}}
\newcommand{\R}{\mathbb{R}}
\newcommand{\C}{\mathbb{C}}
\newcommand{\F}{\mathbb{F}}
\newcommand{\A}{\mathbb{A}}
\newcommand{\PP}{\mathbb{P}}
\newcommand{\E}{\mathbb{E}}
\newcommand{\eps}{\varepsilon}
\newcommand{\dd}{\,\mathrm d}
\newcommand{\sref}[2]{\hyperlink{source:#1:#2}{[S#2]}}
\newif\ifshowcatalogue
\showcataloguetrue
\newcommand{\cataloguescope}[1]{\ifshowcatalogue
 \paragraph{Statement recorded in the supplied source.}
 #1\fi}
\begin{document}
% UnsolvedMath ID 1253; source catalogue code NT-046

\setcounter{section}{1252}

\section{Amicable Integers of Opposite Parity}\label{1253}

{\small\sffamily UnsolvedMath ID 1253 · Number Theory}

\subsection{Definitions and mathematical statement}

\paragraph{Shared notation.}Unless a section says otherwise, $\N=\{1,2,\ldots\}$,
$\Z$ is the ring of integers, $\Q,\R,\C$ are the rational, real, and complex
fields, $[N]=\{1,\ldots,\lfloor N\rfloor\}$, and $\log$ is the natural
logarithm. For real functions of a parameter tending to infinity,
$f\ll g$ and $f=O(g)$ mean $|f|\leq Cg$ eventually for some fixed $C$;
$f\gg g$ reverses this comparison; $f=o(g)$ means $f/g\to0$;
$f\sim g$ means $f/g\to1$; and $f\asymp g$ means both order bounds.
A subscript on $O$ or $\ll$ specifies allowed constant dependence.
The expression $n^{o(1)}$ in an upper bound means growth smaller than
$n^\epsilon$ for each fixed $\epsilon>0$. Local definitions govern any
exception, finite-field convention, or use of the same symbol for another
object. An asymptotic claim is not automatically an effective algorithm.



Positive integers $a\ne b$ are amicable if $\sigma(a)=\sigma(b)=a+b$. Thus the proper positive divisors of each sum to the other. The additional condition is $a\not\equiv b\pmod2$. Equality $a=b$, which would describe a perfect number, is excluded. \sref{1253}{1} \sref{1253}{2}

\paragraph{Statement recorded in the supplied source.}
Do any pairs of amicable numbers exist where one is odd and one is even? \sref{1253}{1}

\paragraph{Residual task reported by the archive.}

The embedded review (2026-08-17) records: Produce an even-odd amicable pair or eliminate all pairs satisfying the known square-form constraints. \sref{1253}{1}

\subsection{Short English statement}

Can an amicable pair consist of one odd integer and one even integer?

\subsection{Research attempt}

\paragraph{Scope and source check.}
The members must be different and of opposite parity. D'Abramo's
primary preprint [S2] states the square-form restrictions and notes
their earlier appearance in work of Gioia. Its abstract was inspected
1 October 2026. The reductions below are proved directly and are not
claimed to be new. Their purpose is to identify the simultaneous
conditions still needed after parity has been used.

\paragraph{Both divisor sums are odd.}
Let the even member be $a$ and the odd member $b$. Since $a+b$ is odd,
both $\sigma(a)$ and $\sigma(b)$ are odd. The prime-power formula
shows that $\sigma(p^e)$ for an odd prime is odd exactly when $e$
is even, while $\sigma(2^r)$ is always odd. Therefore
\[
 a=2^r u^2,\qquad b=v^2,\qquad r\ge1,\quad u,v\text{ odd}.
\]
The two amicability equations now read
\[
 (2^{r+1}-1)\sigma(u^2)=2^r u^2+v^2,
 \qquad \sigma(v^2)=2^r u^2+v^2.
\]
The first determines a candidate square $v^2$ from $r,u$, but the
second is still independent. Omitting that second test would confuse
one direction of the proper-divisor relation with an amicable pair.

\paragraph{The even member cannot be a power of two.}
If $u=1$, the first equation forces $v^2=2^r-1$. For $r=2$ this is
three modulo four, impossible; for $r\ge3$ it is seven modulo eight,
also impossible. For $r=1$, it gives $v=1$, but $(a,b)=(2,1)$ is not
amicable since the proper divisors of one sum to zero. Thus $u>1$ in
every possible opposite-parity pair. The odd member cannot be one
either, because $\sigma(1)=1<a+1$.

This infinite exclusion is stronger than merely checking a list of
even members. It does not exclude even members with a nontrivial odd
square part, and no assumption that such a part has few prime factors
is justified by the parity argument.

\paragraph{A complete one-variable candidate map.}
Define $s(t)=\sigma(t)-t$. Starting with a positive odd $v$, the only
possible even partner of $v^2$ is
\[
                         a=s(v^2).
\]
It is automatically even, because $\sigma(v^2)$ and $v^2$ are odd.
The pair is amicable exactly when $s(a)=v^2$ and $a\ne v^2$.
The latter is automatic from parity. One may first filter by requiring
the odd part of $a$ to be a square, but that filter is necessary rather
than sufficient. This yields an exhaustive search by odd member without
a guess about which member is larger.

Conversely starting from $a=2^r u^2$, put
\[
 b=(2^{r+1}-1)\sigma(u^2)-2^r u^2.
\]
The candidate must be a positive odd square and must satisfy
$\sigma(b)=a+b$. Both formulations are exact and finite at fixed
input. Neither supplies a bound on the input for the existence problem.

\paragraph{Abundancy links the two sizes.}
Write $t=b/a>0$. Then
\[
 I(a)=\frac{\sigma(a)}a=1+t,\qquad
 I(b)=\frac{\sigma(b)}b=1+t^{-1},
 \qquad (I(a)-1)(I(b)-1)=1.
\]
Thus the smaller member is abundant and the larger is deficient;
they cannot both be abundant or both deficient. This conclusion
depends on size, not parity: the equation has not established which
parity belongs to the smaller member. An argument assigning abundance
to the even member without comparison would add an unproved assumption.

For each fixed factor support, the bounds
$I(\prod p^{e_p})<\prod p/(p-1)$ give restrictions on $t$.
But opposite parity does not bound either support, and the two equations
couple their abundancies through a reciprocal relation. Estimating just
one of them from above cannot establish that the exact reciprocal value
is impossible for every other square.

\paragraph{Why scaling known amicable pairs does not change parity.}
If $\{a,b\}$ is amicable and $c$ is coprime to $ab$, multiplicativity
would make $\{ca,cb\}$ amicable only if
$\sigma(c)(a+b)=c(a+b)$, or $\sigma(c)=c$. For $c>1$ the divisors
$1,c$ already show $\sigma(c)>c$, so this fails. Without the
coprimality hypothesis there is no multiplicative identity of this
form at all. In either case multiplying both members by the same
integer cannot turn a same-parity pair into an opposite-parity pair.
Hence a large catalogue of ordinary amicable examples is not itself
a construction for this target.

\paragraph{Executed finite classification.}
The script computes exact divisor sums through $200000$ and tests
every $a<b$ with both members at most that bound. It finds twenty
amicable pairs, none of opposite parity. Its coverage explicitly
requires both members within the bound; a partner larger than it is
not classified by this calculation. Code and result are
\texttt{checks\_second.py}, \texttt{results\_second.json} in
\path{_Local/Erdos_problems/UnsolvedMath/attempt_audit_2026_10_01/number_theory/}.

\paragraph{Remaining gap and outcome.}
The square shapes, exact candidate maps, reciprocal abundancy equation,
and exclusion of a power-of-two member are established. No argument
excludes all nontrivial square parts, and no candidate passing both
proper-divisor equations was found. \textbf{Outcome: UNRESOLVED.}
The finite search and necessary square shapes do not prove nonexistence.

\subsection{Sources}

{\small

\begin{itemize}

\item[\hypertarget{source:1253:1}{S1}] ULAM.AI, \emph{UnsolvedMath}, supplied \texttt{problems.json}, record 1253 (NT-046), statement and background. The embedded literature-review date is 2026-08-17; that is the archive’s review date, not a fresh certification. Dataset landing page: \url{https://huggingface.co/datasets/ulamai/UnsolvedMath}.

\item[\hypertarget{source:1253:2}{S2}] Germano D'Abramo, On Amicable Numbers With Different Parity, arXiv:math/0501402. \url{https://arxiv.org/abs/math/0501402}. Bibliographic information imported from the supplied record.

\item[\hypertarget{source:1253:3}{S3}] Yuta Suzuki, On even-odd amicable pairs, RIMS Kokyuroku 2162 (2020), 127-138. \url{http://hdl.handle.net/2433/261423}. Bibliographic information imported from the supplied record.


\item[E1] Germano D'Abramo, \emph{On Amicable Numbers With Different
Parity}, arXiv:math/0501402v4 (2007).
\url{https://arxiv.org/abs/math/0501402v4}. Abstract and historical
qualification inspected 1 October 2026; the parity reduction is reproved above.

\end{itemize}

}

\label{end:1253}
\end{document}
